\documentclass[10pt]{article}

\usepackage[utf8]{inputenc}

\usepackage[T1]{fontenc}

\usepackage{amsmath}

\usepackage{amsfonts}

\usepackage{amssymb}

\usepackage[version=4]{mhchem}

\usepackage{stmaryrd}

\usepackage{bbold}



\begin{document}

ведений, 5 изд., М., 1971; [14] П р у д н и к о в А. П., Б р ы чк о в Ю. А., М а ри ч е в О. И., Интегралы и ряды. Элементарные функции, М., 1981 ; [15] и х ж е, Интегралы и ряды Специальные функции, М., 1983. Ю.А.Брычков, А.П.Прудников.



СПЕЦИАЛЬНЫІИ АВТОМОРФИЗМ, построенный по автоморфизму $S$ пространства $с$ мерой $(X, v)$ и функции $f$ (заданной на $X$ и принимающей положительные целочисленные значения), - автоморфизм $T$ нек-рого нового пространства с мерой $(M, \mu)$, строящийся следующим образом. Точки $M$ суть пары $(x, n)$, где $x \in X$ и $n$ целое, $0 \leqslant n<f(x)$; при этом $M$ снабжается очевидной мерой $\mu:$ если $A \subset X$ и $f(x)>n$ при всех $x \in A$, то $\mu(A \times$ $\times n)=v(A)$. Если $\mu(M)=\int_{\boldsymbol{X}} f d v<\infty$, то эту меру обычно еще нормируют. Преобразование $T$ увеличивает вторую координату точки $(x, n)$ на единицу, если $n+1<f(x)$, т. е. если при әтом точка не выходит из $M$; в противном случае $T(x, n)=(S x, 0)$. Преобразование $T$ оказывается автоморфизмом пространства с мерой $(M, \mu)$.



Описанная конструкция часто применяется в әргодической теории при построении различных примеров. С другой стороны, роль С. а. ясна из следующего. Отождествляя точку $x \in X \subset(x, 0)$, можно считать, что $X \subset M$. Тогда $f(x)$ есть время, за к-рое точка, находившаяся вначале в $X$ и двигающаяся под действием каскада $\left\{T^{n}\right\}$, возврашается снова в $X$, а $S$ есть производныц̆ автоморфизм $T_{X}$. Таким образом, с помощью С. а. как бы восстанавливают траектрии динамич. системы во всем фазовом пространстве, наблюдая только прохождения движущейся точки через множество $X$.



СПЕШИАЛЬНЫИ ПОТОК, построенный По В.Аносов. физму $S$ пространства с мерой $(X, v)$ и измеримой функции $f$, заданной на $X$ и принимающей положительные значения, - измеримый поток в нек-ром новом пространстве с мерой $(M, \mu)$, строящийся следующим образом. Точки $M$ суть пары $(x, s)$, где $x \in X$ и $0 \leqslant s<f(x)$; мера $\mu$ - ограничение на множестве $M$, рассматриваемом как подмножество в $X \times[0, \infty)$ - прямого іроизведения меры $v$ на $X$ и меры Лебега на $[0, \infty)$. Если $\mu(M)=\int_{X} f d v<\infty$, то эту меру обычно еще нормируют. Движение же происходит таким образом, что вторая координата точки $(x, s)$ увеличивается с единичной скоростью, пока не достигает значения $f(x)$; а әтот момент времени точка перескакивает в положение $(S x, 0)$.



В әргодической теории С. П. играет роль, аналогичную роли сечений и последования отображений при исследовании гладких динамич. систем: $X$ играет роль сечения, а $S$ - отображения последования. Но в топологич. теории построить сечение в виде многообразия удается, вообще говоря, только локально. В эргодической же теории построение глобального сечения возможно при очень широких условиях, ибо здесь нет ограничений, связанных с топологией. (Если даже исходный поток является гладким, все равно допускается, чтобы секущая поверхность была разрывной.) Поэтому при очень пироких условиях измеримый поток метрически изоморфен нек-рому С. п., даже с нек-рыми дополнительными условиями на $f$ (см. [1]). Родственным понятием является специиальныц̆ автоморфизм.



Лum.: [1] К о р н ф е ль И. П., С ин а й Я. Г., Ф ом и н С. В., Әргодическая теория, М., 1980. Д.В.'Аносов. СПИН - одна из величин, характеризующая внутренние степени свободы квантовой частицы (или квантового поля). Нерелятивистская частица имеет спин $S$ $\left(S=0, \frac{1}{2}, 1, \frac{3}{2}, 2, \ldots\right)$, если ее волновая функция $\psi(x)$ принимает значения из нек-рого линейного пространства $L$, в к-ром действует неприводимое (однозначное или двузначное) представление $g \rightarrow T g$ веса $S$ групиы вращений $\mathrm{SO}_{3}$ трехмерного пространства (см. [1], [2]) так, что ири переходе от одной ортогональной системы координат в $\mathbb{R}$ к другой при помощи матрицы $g \in \mathrm{SO}_{3}$ значения волновой функции преобразуются матрицей $T g$. В случае релятивистской частицы, волновая функция к-рой преобразуется с помощью какогонибудь представления $g \rightarrow T g$ группы Јоренца, ее С. наз. максимальный из весов неприводимых представлений группы вращений $S O_{3}$, на к-рые раскладывается представление $g \rightarrow T g$.



Лum.: [1] Л а нд а у Л. Д., Л и ф ш и ц Е. М., Квантовая механика, 3 изд., М., 1974; [2] Г е л ь ф а н д И. М., М и нл о с Р. А., II а П и р о 3. Я., Представления группы вращений и группы Лоренца, их применения, М., 1958.



P. А. Минлос СПИНОР - әлемент пространства спинорного представления. Напр., если $Q$ - невырожденная квадратичная форма в $n$-мерном пространстве $V$ над полем $k$, имеющая максимальный индекс Витта $m=[n / 2]$ (последнее условие всегда выполнено, если поле $k$ алгебраически замкнуто), то в качестве пространства С., отвечающего форме $Q$, можно рассматривать внешнюю алгебру над максимальными (размерности $m$ ) вполне изотропными подпространствами в $V$.



С. были рассмотрены впервые в 1913 Э. Картаном (E. Cartan) в его исследованиях по теории представлений непрерывных групп и вновь открыты в 1929 Б. Ј. Ван дер Варденом (B. L. van der Waerden) в исследованиях по квантовой механике [так, было обнаружено, что явление спина электрона и других элементарных частиц описывается физич. величинами, не принадлежащими к известным ранее типам величин (тензорам, псевдотензорам и т. Д.), напр., они определяются липь с точностью до знака, и при повороте системы координат на угол $2 \pi$ вокруг нек-рой оси все компоненты әтих величин меняют знак].



В настоящее время спинорное исчисление нашло широкое применение во многих отраслях математики и позволило решить ряд трудных задач алгебраической и диффференциальной топологии (напр., задача о числе ненулевых векторных полей на $k$-мерной сфере, задача об индексе эллиптического оператора, $K$-теория).



Лum.: [1] д у б р о в и н Б. А., Н о в и к о в С. П., Ф ом е н к о А. Т.' Современная геометрия, М., 1979; [2] Ћ е ли механике, М., $1982 ;$ [3] К а р т а н Ә., Теория спиноров, пер. с франц., М., 1947; [4] К а р у б и М., К-теория. Введение, пер. с форанц., М., 1947; [4] К а р у 6 и М., К-теория. Введение, пер.

с англ., М., 1981.

М. И. Войеховский.



СПЙНОЬРНАЯ ГРУППА невырожденной квадратитной формы $Q$ на $n$-мерном $(n \geqslant 3)$ векторном пространстве $V$ над полем $k$ - связная линейная алгебраич. группа, являющаяся универсальной накрывающей неприводимой компоненты единицы $O_{n}^{+}(Q)$ ортогональной групшы $O_{n}(Q)$ формы $Q$. Если char $k \neq 2$, то группа $O_{n}^{+}(Q)$ совпадает со специальной ортогональной групиой $S O_{n}(Q)$. С. г. строится следующим образом. Пусть $C=$ $=C(Q)$ - Клифбфорда алгебра пары $(V, Q), C+$ (соответственно $C^{-)}$- подпространство в $C$, порожденное произведениями четного (соответственно нечетного) числа элементов из $V, \beta$ - канонич. антиавтоморфизм алгебры $C$, определяемый формулой



\[

\beta\left(v_{1} v_{2} \ldots v_{k}\right)=v_{k} \ldots v_{2} v_{1} .

\]



Вложение $V \subset C$ позволяет определить $\mathbf{r} \mathbf{p} \mathbf{y}$ п п у I л и фф о о д а



\[

G=\left\{s \in C \mid s \text { обратим в } C \text { п } s V s^{-1}=V\right\}

\]



и ч е тн ую (или спе пи а льн ую) $г \mathbf{p}$ у п п у Кли фф о р д а



\[

G^{+}=G \bigcap G^{+}

\]



C. г. $\operatorname{Spin}_{n}=\operatorname{Spin}_{n}(Q)$ определяется равенством



\[

\operatorname{Spin}_{n}=\{s \in G+\mid s \beta(s)=1\} \text {. }

\]



C. г. $\operatorname{Spin}_{n}-$ простая (при $n \neq 4$ ) связная односвязная линейная алгебраич. группа типа $B_{m}$ при $n=2 m+1$





\end{document}